\chapter{Generative Models and Synthetic Data}\label{ch:ml:generative-models-and-synthetic-data}

A \omterm{def:ml:generative-models-and-synthetic-data:gan}{generative adversarial network} trained on twenty years of daily index returns produces windows on which a simple
strategy, trained and tested on generated data, earns a Sharpe ratio of 10.7. Trained on the same generated data and
tested on the ten real years that followed, it earns 0.5. The network had invented a predictability the market did not
have, and it had not learned the one the market did. Synthetic data are one of machine learning's most useful tools in
finance: they let a desk stress a strategy on crashes that happened once, test a backtester on paths with a known truth,
and train models where data are scarce. They also mislead in ways that no stylised-fact checklist catches. This chapter
builds five generators, from a block bootstrap to a \omterm{def:ml:generative-models-and-synthetic-data:diffusion}{diffusion model}, and judges them by four tests; its data are
synthetic, so that the real series has a known structure to learn.

\section{What a generative model is for}

\begin{definition}[Generative model]\label{def:ml:generative-models-and-synthetic-data:gen}
A \emph{generative model}\index{generative model} is a model of the distribution of the data that can draw new samples from
it. For market data the samples are paths (of returns, order-book states, curves) or windows of them, and the model is judged
by whether the samples are useful for the purpose, not only by whether they look real.
\end{definition}

The chapter's real series is thirty years of daily index returns from a GJR-GARCH model with Student-t shocks (Book~4,
chapter~18): fat tails, volatility clustering, a leverage effect, six planted crashes of 8 to 11 standard deviations, and
one piece of predictability, a conditional mean of 0.1 times the current volatility in the direction of the last twenty
days' return. The first twenty years train; the last ten test. Every generator produces windows of 32 trading days. Its
purposes set the tests: stress testing needs the tails and the clustering; training a predictor needs the predictability,
and needs it not to be invented; privacy or data sharing needs samples that are not copies.

Two classical generators set the baseline. A block bootstrap (Book~4, chapter~13) cuts the training series into 20-day
blocks and glues random blocks together: it keeps everything inside a block and nothing across blocks. A GARCH(1,1) model
with Student-t shocks, fitted by maximum likelihood ($\alpha = 0.079$, $\beta = 0.886$, 3.6 degrees of freedom), keeps
what its equations say and nothing else: it has no leverage effect and no momentum.

\section{Variational autoencoders and adversarial networks}

\begin{definition}[Variational autoencoder]\label{def:ml:generative-models-and-synthetic-data:vae}
A \emph{variational autoencoder}\index{variational autoencoder} (VAE) is an \omterm{def:ml:cross-sectional-deep-models:ae}{autoencoder} (chapter~9) whose encoder outputs a
distribution over a low-dimensional latent code and whose decoder outputs a distribution over the data; it is trained to
maximise a lower bound on the likelihood, the reconstruction log-likelihood minus the Kullback--Leibler divergence of the
encoder's distribution from a standard normal prior, and it generates by decoding codes drawn from the prior (Kingma and
Welling, 2014).
\end{definition}

\begin{definition}[Generative adversarial network, mode collapse]\label{def:ml:generative-models-and-synthetic-data:gan}
A \emph{generative adversarial network}\index{generative adversarial network} (GAN) trains two networks against each other:
a generator that maps noise to samples, and a discriminator that tells generated samples from real ones; the generator is
trained to make the discriminator fail (Goodfellow and co-authors, 2014). \emph{Mode collapse}\index{mode collapse} is the
failure in which the generator covers only part of the data's distribution (a few typical patterns, no extremes) because
that is enough to fool the current discriminator.
\end{definition}

The chapter's VAE has a four-dimensional latent code and outputs a mean and a variance for each of the 32 days; its GAN
maps 16-dimensional noise through two layers of 128 units. Both are \omterm{def:ml:neural-networks-for-noisy-tabular-data:mlp}{multilayer perceptrons} on windows standardised by the
training returns' volatility, trained for 4\,000 steps of \omterm{def:ml:neural-networks-for-noisy-tabular-data:adam}{Adam} on one core in a few seconds. The VAE's windows have excess
kurtosis 11.1 against the real training windows' 13.0; the GAN's have 2.0: it produces calm and moderately volatile weeks
and never a crash, the signature of \omterm{def:ml:generative-models-and-synthetic-data:gan}{mode collapse} on a fat-tailed distribution. Quant GANs (Wiese and co-authors, 2020) and
the small-data market simulator of Buehler and co-authors (2020) are designs aimed at these failures.

\section{Diffusion models}

\begin{definition}[Diffusion model]\label{def:ml:generative-models-and-synthetic-data:diffusion}
A \emph{diffusion model}\index{diffusion model} learns to reverse a process that gradually adds Gaussian noise to the data:
a network is trained to predict the noise added to a sample at each of $T$ steps, and new samples are generated by starting
from pure noise and removing the predicted noise step by step (Ho, Jain and Abbeel, 2020).
\end{definition}

\Cref{lst:ml:gen:diffusion} is a complete denoising \omterm{def:ml:generative-models-and-synthetic-data:diffusion}{diffusion model} on windows: 100 noise levels, a two-layer network that
takes the noisy window and the step, and the reverse chain. It is the most stable of the three networks to train, and on
this data the weakest at the tails (kurtosis 4.0): the Gaussian noise it starts from and the few thousand training steps
leave the extremes under-represented.

\section{Judging synthetic data}

\begin{definition}[Classifier two-sample test, train-on-synthetic test-on-real]\label{def:ml:generative-models-and-synthetic-data:tests}
The \emph{classifier two-sample test}\index{classifier two-sample test} trains a classifier to tell generated samples from
real ones and reports its accuracy on held-out samples: 50\% means the classifier cannot tell them apart (Lopez-Paz and
Oquab, 2017). \emph{Train-on-synthetic test-on-real}\index{train-on-synthetic test-on-real} (TSTR) trains a downstream model
on generated data and evaluates it on real data, comparing with the same model trained on real data (Esteban, Hyland and
Rätsch, 2017).
\end{definition}

Four tests judge the five generators (\cref{tab:ml:gen:judge}). The \emph{scorecard} checks six stylised facts (Cont,
2001) within windows: heavy tails (excess kurtosis above 5), no linear autocorrelation, volatility clustering in absolute
returns at lags 1 and 10, the leverage effect (negative correlation of a return with the next absolute return), gain-loss
asymmetry (negative skewness) and aggregational Gaussianity (kurtosis of 20-day sums below half the daily one). The real
training windows pass all six. The \emph{two-sample test} uses a gradient-boosted classifier on 13 window features; real
windows overlap, so its halves are split by year, or the classifier recognises a test window by its neighbours in the
training half (a random split scores the even training years against the odd ones at 0.83). \emph{TSTR} fits a linear
forecast of a window's last day from features of its first 31 on 20\,000 generated windows and trades the sign of the
forecast on the ten real test years; it is averaged over five generated samples. \emph{Train-on-synthetic
test-on-synthetic} (TSTS) does the same with a second generated sample as the \omterm{def:ml:why-financial-machine-learning-is-different:sets}{test set}. The last column is the
\emph{\omterm{def:ml:large-language-models-in-finance:memo}{memorisation} check} (\cref{lst:ml:gen:tests}): the share of generated windows closer to a training window than 95\% of
the real test windows are.

\begin{table}[htbp]
\centering\small
\begin{tabular}{@{}lcccccc@{}}
\toprule
& facts & two-sample & \multicolumn{2}{c}{TSTR} & TSTS & closer than\\
\cmidrule(lr){4-5}
generator & passed & accuracy & IC & Sharpe & Sharpe & real (\%)\\
\midrule
real (train on real, test on real) & 6 & 0.545 & 0.078 & 0.86 & & 5.0\\
block bootstrap & 6 & 0.491 & 0.062 & 0.71 & 1.19 & 14.5\\
GARCH-t & 4 & 0.539 & $-0.019$ & $-0.10$ & $-0.09$ & 5.7\\
VAE & 6 & 0.619 & 0.040 & 0.76 & 0.79 & 5.1\\
GAN & 3 & 0.749 & 0.052 & 0.52 & 10.68 & 11.2\\
diffusion & 4 & 0.693 & 0.030 & 0.23 & 1.45 & 1.4\\
\bottomrule
\end{tabular}
\caption{Five generators judged four ways. Two-sample accuracy of the real row: the training years' even against odd
years. Data: \texttt{ml\_gen.judge}.}\label{tab:ml:gen:judge}
\end{table}

\begin{omfigure}
\begin{tikzpicture}
\begin{axis}[width=11.5cm,height=6cm,xlabel={lag (days)},ylabel={autocorrelation of $|r_t|$},tick label style={font=\small},
  label style={font=\small},xmin=1,xmax=20,ymin=0,ymax=0.32,grid=major,grid style={gray!25},legend columns=3,
  legend style={font=\footnotesize,at={(0.5,-0.25)},anchor=north,draw=none,fill=none,
  /tikz/every even column/.append style={column sep=6pt}},table/col sep=comma]
\addplot[black,ultra thick] table[x=lag,y=real]{figdata/ml/16-generative-models-and-synthetic-data/absacf.csv};
\addplot[omDef,thick,mark=*,mark size=1.3pt] table[x=lag,y=boot]{figdata/ml/16-generative-models-and-synthetic-data/absacf.csv};
\addplot[gray,thick,mark=triangle*,mark size=1.3pt] table[x=lag,y=garch]{figdata/ml/16-generative-models-and-synthetic-data/absacf.csv};
\addplot[omProp,thick,mark=square*,mark size=1.3pt] table[x=lag,y=vae]{figdata/ml/16-generative-models-and-synthetic-data/absacf.csv};
\addplot[omThm,thick,mark=diamond*,mark size=1.5pt] table[x=lag,y=gan]{figdata/ml/16-generative-models-and-synthetic-data/absacf.csv};
\addplot[omProp!50,thick,dashed] table[x=lag,y=diff]{figdata/ml/16-generative-models-and-synthetic-data/absacf.csv};
\legend{real,block bootstrap,GARCH-t,VAE,GAN,diffusion}
\end{axis}
\end{tikzpicture}

\omcaption{Volatility clustering: autocorrelation of absolute daily returns within 32-day windows, real training windows
against 4\,000 windows from each generator. Data: \texttt{ml\_gen.abs\_acf}.}\label{fig:ml:gen:acf}
\end{omfigure}

The table holds the chapter's lessons. The scorecard is the weakest test: the VAE passes all six facts and is told apart
from real windows 62\% of the time; the block bootstrap passes all six and is indistinguishable (0.491) because its windows
are pieces of real ones. The two-sample test sees what the scorecard misses, but not what matters for a strategy. Only TSTR
measures that: the real model earns 0.86; the bootstrap's windows keep most of the momentum inside their 20-day blocks
(0.71); the VAE learned a noisy version of it (0.76, with a standard deviation of 0.28 across samples); GARCH, which cannot
represent it, teaches nothing ($-0.10$, standard deviation 0.65: pure noise); the \omterm{def:ml:generative-models-and-synthetic-data:diffusion}{diffusion model} teaches little (0.23).

\section{The ways synthetic data mislead}

\Cref{fig:ml:gen:sharpe} is the hook's result. On the GAN's own windows a strategy earns a Sharpe ratio of 10.7: the
generator produces far more linear autocorrelation (0.099 at lag 1, against 0.078 in the real windows) and persistent
volatility patterns that a linear model can trade, an artefact of a network that has collapsed onto a few smooth shapes.
Tested on real data, the same strategy earns 0.52. Every generator except GARCH shows a gap between TSTS and TSTR, and
nothing inside the synthetic world reveals it.

\begin{omfigure}
\begin{tikzpicture}
\begin{axis}[width=11cm,height=6cm,ybar,bar width=10pt,ymin=-1,ymax=11.5,ylabel={Sharpe ratio (annualised)},area legend,
  xtick={0,1,2,3,4},xticklabels={bootstrap,GARCH-t,VAE,GAN,diffusion},tick label style={font=\small},
  label style={font=\small},xmin=-0.5,xmax=4.5,grid=major,grid style={gray!25},legend columns=3,
  legend style={font=\footnotesize,at={(0.5,-0.2)},anchor=north,draw=none,fill=none},table/col sep=comma]
\addplot[fill=omThm,draw=none,bar shift=-5.5pt] table[x=i,y=tsts]{figdata/ml/16-generative-models-and-synthetic-data/sharpe.csv};
\addplot[fill=omDef,draw=none,bar shift=5.5pt] table[x=i,y=tstr]{figdata/ml/16-generative-models-and-synthetic-data/sharpe.csv};
\addplot[black!70,dashed,thick,sharp plot,forget plot] table[x=x,y=y]{figdata/ml/16-generative-models-and-synthetic-data/trtr.csv};
\legend{tested on generated windows,tested on real windows}
\end{axis}
\end{tikzpicture}

\omcaption{The same strategy trained on each generator's windows, tested on a second generated sample and on the ten real
test years; dashed, trained and tested on real data. Data: \texttt{ml\_gen.judge}.}\label{fig:ml:gen:sharpe}
\end{omfigure}

The \omterm{def:ml:large-language-models-in-finance:memo}{memorisation} column shows the other two failures. The block bootstrap's windows are closer to a training window than
the real test windows are three times as often as chance (14.5\% against 5\%): they are copies by construction, which
is fine for stress tests and useless where synthetic data must protect the real data's owner. The GAN's 11.2\% has a
different cause: its windows are all calm, and calm windows have close neighbours among the many calm training windows.
Nearest-neighbour distance flags copying and collapse alike, and has to be read with the tails. The \omterm{def:ml:large-language-models-in-finance:memo}{memorisation} of
chapter~14 applies to generators as to \omterm{def:ml:large-language-models-in-finance:lm}{language models}: a large generator trained long enough on a small market history
reproduces it.

\begin{method}[Using synthetic market data]\label{met:ml:gen:method}
\begin{enumerate}
\item Name the purpose (stress, training, testing a backtester, sharing), and choose the tests it needs.
\item Start from the classical generators (bootstrap, GARCH, a simulator with known mechanisms) and require a network to
beat them on those tests.
\item Judge by TSTR against train-on-real, never by results inside the synthetic world; add the two-sample test with a
blocked split and the \omterm{def:ml:large-language-models-in-finance:memo}{memorisation} check.
\item Treat any edge found only in synthetic data as an artefact until it is found in real data.
\end{enumerate}
\end{method}

\section{Tutorial: paths that never happened}\label{tut:ml:generative-models-and-synthetic-data:tut}

\begin{tutorial}
\textbf{Goal.} Fit five generators to the training years, and judge them by scorecard, two-sample test, TSTR, TSTS and
nearest-neighbour distance. \textbf{End state:} \cref{tab:ml:gen:judge}, \cref{fig:ml:gen:acf,fig:ml:gen:sharpe}.
\begin{tutsteps}
\item \textbf{A \omterm{def:ml:generative-models-and-synthetic-data:gan}{generative adversarial network}.}
\omcode{code/firm/genmkt/firm_genmkt.py}{166}{195}{A generator and a discriminator trained against each other.}\label{lst:ml:gen:gan}
\item \textbf{A denoising \omterm{def:ml:generative-models-and-synthetic-data:diffusion}{diffusion model}.}
\omcode{code/firm/genmkt/firm_genmkt.py}{198}{235}{Noise, a denoiser, and the reverse chain.}\label{lst:ml:gen:diffusion}
\item \textbf{Train on synthetic, test on real.}
\omcode{code/firm/genmkt/firm_genmkt.py}{308}{319}{The TSTR score: rank IC and Sharpe ratio on real windows.}\label{lst:ml:gen:tests}
\item \textbf{Run} \texttt{ml\_gen.judge()} (about a minute on one core) and \texttt{fig\_gen.py}.
\end{tutsteps}
\textbf{What to change next.} Condition the generators on the volatility regime; train the GAN longer and watch its TSTS
Sharpe ratio and its two-sample accuracy; replace the block bootstrap's 20-day blocks by 5-day blocks and watch the
momentum disappear from its TSTR.
\end{tutorial}

\section{Build: generators and their judges}\label{bld:ml:generative-models-and-synthetic-data:genmkt}

\begin{build}
\bfield{Purpose} Synthetic market data with a stated purpose, judged by tests that match it.
\bfield{Interface} \texttt{real\_index(n\_days, seed)}, \texttt{windows}; \texttt{block\_bootstrap}, \texttt{fit\_garch},
\texttt{garch\_windows}; \texttt{VAE}, \texttt{GAN}, \texttt{Diffusion} with \texttt{fit(W, steps, seed)} and
\texttt{sample(n, seed)}; \texttt{facts}, \texttt{scorecard}, \texttt{c2st(A, B, groups\_a, groups\_b, seed)},
\texttt{tstr(W\_train, W\_test)}, \texttt{nn\_distance}.
\bfield{Rules} Every generator is fitted on the training period only; overlapping real windows are split by time; TSTR is
averaged over samples; results inside the synthetic world are never reported alone.
\bfield{Acceptance tests} \texttt{code/firm/genmkt/tests/}: the GARCH fit recovers planted parameters; the bootstrap copies
blocks; the two-sample test is near 0.5 on two samples of one distribution and near 1 on two different ones; TSTR recovers a
planted predictor; the networks train deterministically and reduce their losses; nearest-neighbour distance is zero for
copies.
\bfield{Stretch} Conditional generation on regimes; signature-based generators; Book~10's agent-based market (chapter~27)
as a generator with mechanisms, judged by the same tests.
\end{build}

\begin{omsources}
\item D.~P.~Kingma and M.~Welling, ``Auto-encoding variational Bayes'', arXiv:1312.6114, 2014.
\item I.~Goodfellow and co-authors, ``Generative adversarial nets'', arXiv:1406.2661, 2014.
\item J.~Ho, A.~Jain and P.~Abbeel, ``Denoising diffusion probabilistic models'', arXiv:2006.11239, 2020.
\item M.~Wiese, R.~Knobloch, R.~Korn and P.~Kretschmer, ``Quant GANs: deep generation of financial time series'',
\emph{Quantitative Finance} 20(9), 2020.
\item H.~Buehler, B.~Horvath, T.~Lyons, I.~Perez Arribas and B.~Wood, ``A data-driven market simulator for small data
environments'', arXiv:2006.14498, 2020.
\item D.~Lopez-Paz and M.~Oquab, ``Revisiting classifier two-sample tests'', arXiv:1610.06545, 2017.
\item R.~Cont, ``Empirical properties of asset returns: stylized facts and statistical issues'', \emph{Quantitative Finance}
1(2), 2001.
\item C.~Esteban, S.~L.~Hyland and G.~Rätsch, ``Real-valued (medical) time series generation with recurrent conditional
GANs'', arXiv:1706.02633, 2017.
\end{omsources}

\section{Exercises}

\begin{exercise}[$\star$]\label{exo:ml:generative-models-and-synthetic-data:1}
A two-sample test scores 0.62 on 4\,000 generated against 4\,000 real windows (2\,000 of each held out). Is the difference
from 0.5 significant? What would 0.51 tell you?
\end{exercise}

\begin{exercise}[$\star$]\label{exo:ml:generative-models-and-synthetic-data:2}
Which of the six stylised facts can a GARCH(1,1) model with symmetric Student-t shocks never produce, and why?
\end{exercise}

\begin{exercise}[$\star$]\label{exo:ml:generative-models-and-synthetic-data:3}
The fitted GARCH has $\alpha + \beta = 0.964$. What is the half-life of a volatility shock in days?
\end{exercise}

\begin{exercise}[$\star\star$]\label{exo:ml:generative-models-and-synthetic-data:4}
Why does the block bootstrap keep the momentum effect partly, and what block length would destroy it?
\end{exercise}

\begin{exercise}[$\star\star$]\label{exo:ml:generative-models-and-synthetic-data:5}
Why did a random split make the two-sample test score real windows against real windows at 0.83?
\end{exercise}

\begin{exercise}[$\star\star$]\label{exo:ml:generative-models-and-synthetic-data:6}
\emph{Find the flaw.} ``Our GAN's paths pass all our stylised-fact tests, and a strategy trained on a million of them has a
Sharpe ratio of 3 out of sample on another million, so it is robust.''
\end{exercise}

\begin{exercise}[$\star\star\star$]\label{exo:ml:generative-models-and-synthetic-data:7}
\emph{Coding.} Train the GAN for 1\,000 and for 8\,000 steps. Report its excess kurtosis, two-sample accuracy and TSTS
Sharpe ratio each time. Does longer training cure the collapse?
\end{exercise}

\begin{exercise}[$\star\star\star$]\label{exo:ml:generative-models-and-synthetic-data:8}
Derive the VAE's objective: show that $\log p(x)\ge\E_{q(z\mid x)}[\log p(x\mid z)] - \mathrm{KL}(q(z\mid x)\,\|\,p(z))$, and
write the KL term for a diagonal Gaussian $q$ and a standard normal prior.
\end{exercise}

\section{Problem: Paths That Never Happened}

\begin{problem}[{Weekend problem --- shape is not predictability}]\label{pb:ml:generative-models-and-synthetic-data:1}
The chapter's thirty years of returns, five generators and four tests.

\medskip\noindent\textbf{Part I --- The data and the generators.}
\begin{enumerate}
\item What structure does the real series have, and which parts can each generator represent?
\item What are the fitted GARCH parameters, and what does the model miss?
\item What do the VAE and GAN networks look like, and how long do they train?
\item What is the \omterm{def:ml:generative-models-and-synthetic-data:diffusion}{diffusion model}'s generating process?
\end{enumerate}

\medskip\noindent\textbf{Part II --- Judging.}
\begin{enumerate}[resume]
\item How many stylised facts does each generator pass?
\item What are the two-sample accuracies, and why must real windows be split by time?
\item What are the TSTR information coefficients and Sharpe ratios against train-on-real?
\item What does the \omterm{def:ml:large-language-models-in-finance:memo}{memorisation} check show?
\end{enumerate}

\medskip\noindent\textbf{Part III --- Misleading.}
\begin{enumerate}[resume]
\item What Sharpe ratio does the GAN's world give, and where does it come from?
\item Why is GARCH's TSTR noise, and how large is that noise?
\item Which generator would you use for a stress test, and which to augment a \omterm{def:ml:why-financial-machine-learning-is-different:sets}{training set}?
\item Why is the bootstrap unsuitable for sharing data outside the firm?
\end{enumerate}

\medskip\noindent\textbf{Part IV --- The verdict.}
\begin{enumerate}[resume]
\item State the \emph{named result}: each generator's scorecard pass count and two-sample accuracy, and the TSTR gap in
Sharpe ratio.
\item What would convince you that a generator has learned a real predictability?
\item How would you test a backtester with a generator?
\item When is a mechanistic simulator better than a learned generator?
\item How would you condition a generator on a regime, and what new risk does it bring?
\item What would you log about a synthetic dataset used in research?
\item What does the \omterm{def:ml:generative-models-and-synthetic-data:diffusion}{diffusion model}'s weakness at the tails imply for stress tests?
\item In one sentence: what does a \omterm{def:ml:generative-models-and-synthetic-data:gen}{generative model} learn from twenty years of returns?
\end{enumerate}
\end{problem}

\section{Interview questions}

\begin{interviewq}[$\star$ \iqroles{researcher, mle}]\label{iq:ml:generative-models-and-synthetic-data:1}
Name the main stylised facts of daily returns. Which does a GARCH model reproduce?
\end{interviewq}

\begin{interviewq}[$\star\star$ \iqroles{mle}]\label{iq:ml:generative-models-and-synthetic-data:2}
Explain the difference between a VAE, a GAN and a \omterm{def:ml:generative-models-and-synthetic-data:diffusion}{diffusion model} in one paragraph each, and one failure of each.
\end{interviewq}

\begin{interviewq}[$\star\star$ \iqroles{researcher}]\label{iq:ml:generative-models-and-synthetic-data:3}
How would you evaluate a generator of synthetic market data?
\end{interviewq}

\begin{interviewq}[$\star\star$ \iqroles{researcher, trader}]\label{iq:ml:generative-models-and-synthetic-data:4}
A colleague augments the \omterm{def:ml:why-financial-machine-learning-is-different:sets}{training set} with GAN paths and the backtest improves. What do you check?
\end{interviewq}

\begin{interviewq}[$\star\star$ \iqroles{mle}]\label{iq:ml:generative-models-and-synthetic-data:5}
What is \omterm{def:ml:generative-models-and-synthetic-data:gan}{mode collapse}, and how would you detect it in generated returns?
\end{interviewq}

\begin{interviewq}[$\star\star\star$ \iqroles{researcher}]\label{iq:ml:generative-models-and-synthetic-data:6}
Design a \omterm{def:ml:generative-models-and-synthetic-data:tests}{classifier two-sample test} for overlapping windows of a time series. What goes wrong with a random split?
\end{interviewq}
